\chapter{Topology}

In this chapter, we introduce essential vocabulary used in \emph{topology}. Each term is accompanied by a 
brief explanation and its formal mathematical definition.

\section{Introduction to topological nomenclature}

\subsection{Tological Space (Informal)}

A \emph{topological space} is a mathematical space whose elements are called points and which obeys certain axioms.
The formal definition will come later, but for the sake of this chapter, we will introduce it later.

\subsection{Open Ball \texorpdfstring{\( B_\varepsilon(x) \)}{}} 

For a metric space \( (X, d) \), the \emph{open ball} centered at \( x \in X \) with radius \( \varepsilon > 0 \) is defined as: 
	      
\[
	B_\varepsilon(x) := \{ y \in X \mid d(x, y) < \varepsilon \} = (x - \varepsilon, x + \varepsilon).
\]

It represents the set of all points within distance \( \varepsilon \) from \(x\), excluding the boundary.

\subsection{Open Set}
	     
A subset \( U \subseteq X \) of a topological space is called \emph{open} if for every point \( x \in U \), there exists an \( \varepsilon > 0 \) 
such that the open ball \( B_\varepsilon(x) \subseteq U \). 

Intuitively, an open set contains none of its boundary points and every point has some 
wiggle room around it.

\textbf{Example:}

We will prove that  \((0,1)\) is an open set. 

This means \(x_0 > 0\) and \(1 - x_0 > 0\). First we define 

\[
	\varepsilon := \frac{1}{2} \min {x_0, 1 - x_0};
\]

then from \(|y - x_0| \le \varepsilon\) we get 

\begin{align*}
	&y \ge x_0 |y - x_0 | \ge x_0 - \varepsilon \ge \frac{1}{2}x_0 > 0, \\ 
	&y \le x_0 |y - x_0 | \le x_0 + \varepsilon \le x_0 + \frac{1}{2}(1 - x_0) = \frac{1}{2} + \frac{1}{2}x_0 < 1.
\end{align*}

\QED 

\subsection{The Interior of a Set} 

The \emph{interior} denoted as \(A^{\circ}\) of a set \(A\) is the set of all inner points of \(A\). 

\subsection{Closed Set} 
	      
A subset \( A \subseteq X \) is called \emph{closed} if its complement \( X \setminus A \) is open. Equivalently, \(A\) contains all its limit points. 
That is, \(A\) is closed if it includes its boundary.

\textbf{Example:} \([0,1]\)

\subsection{Boundary Point}
	      
A point \( x \in X \) is a \emph{boundary point} of a set \( A \subseteq X \) if every open ball around \(x\) contains both points in \(A\) 
and in \( X \setminus A \). The set of all boundary points is called the \emph{boundary} of \(A\), denoted \( \partial A \).

\subsection{Accumulation Point / Limit Point}

A point \( x \in X \) is an \emph{accumulation point} of a set \( A \subseteq X \) if every open ball \( B_\varepsilon(x) \) contains a 
point of \( A \setminus \{x\} \). In other words, points of \(A\) cluster arbitrarily close to 
\(x\), even if \( x \notin A \).

\subsection{Isolated Point} 
	      
A point \( x \in A \) is an \emph{isolated point} if there exists \( \varepsilon > 0 \) such that \( B_\varepsilon(x) \cap A = \{x\} \). 
That is, \(x\) stands alone in \(A\) without other points of \(A\) nearby.

\textbf{Example:}

\(-1\) is an isolated point of \({-1} \cup (0,1)\).

\subsection{Boundary}

The set \(A \cup \{\text{ all boundary points }\}\) is called the \emph{boundary} of \(A\). It is denoted as \(\overline{A}\).

\subsection{Compact Set} 

A set \( K \subseteq X \) is \emph{compact} if every open cover of \( K \) has a finite sub-cover. 
In \(\R^n\), this is equivalent to \( K \) being closed and bounded (by the Heine–Borel theorem).

\subsection{Dense Set} 
	      
A subset \( D \subseteq X \) is \emph{dense} in \(X\) if every point \( x \in X \) is either in \( D \) or is a limit point of \( D \). 
Equivalently, the closure of \( D \) is \(X\), i.e., \( \overline{D} = X \).


